%%% KLOCKI: MIEJSCA NA RYSUNKI
%%% Plik zapisz w folderze figures pod podaną nazwą (pdf, png albo jpg), a rysunek pojawi się sam.
%%% Podpisy są robocze. Wykresy tworzy skrypt Analiza/wykresy.py od razu w folderze figures (python Analiza/wykresy.py);
%%% szerokość pliku jest równa szerokości w nawiasie kwadratowym, więc nie zmieniaj jej bez ponownego uruchomienia skryptu.
%%% Zrzuty ekranu przygotowujesz we własnym zakresie. Opis pokazuje tylko, z czego powstaje rysunek.

\MiejsceNaRysunek{architektura}{Zależności bibliotek kodu}{fig:architektura}{Schemat 4 bibliotek, \cref{tab:biblioteki}.}

\MiejsceNaRysunek{modele}{Modele zestawu kafelków}{fig:modele}{Podgląd 43 modeli; scena z \kol{TilesetPreviewSceneBuilder}.}

\MiejsceNaRysunek{orientacja}{Etykiety ścian i obroty kafelka}{fig:orientacja}{Przykład: \kol{Mur_Dol} w 4 obrotach.}

\begin{figure}
  \centering
  \begin{subfigure}[t]{0.32\linewidth}
    \centering
    \MiejsceNaObraz[\linewidth]{krok-1}{Zrzut z podglądu krok po kroku}
    \caption{Etap pierwszy}
  \end{subfigure}\hfill
  \begin{subfigure}[t]{0.32\linewidth}
    \centering
    \MiejsceNaObraz[\linewidth]{krok-2}{Zrzut z podglądu krok po kroku}
    \caption{Etap drugi}
  \end{subfigure}\hfill
  \begin{subfigure}[t]{0.32\linewidth}
    \centering
    \MiejsceNaObraz[\linewidth]{krok-3}{Zrzut z podglądu krok po kroku}
    \caption{Etap trzeci}
  \end{subfigure}
  \caption{Generacja mapy krok po kroku}
  \label{fig:krok-po-kroku}
\end{figure}

\MiejsceNaRysunek{okno-generatora}{Okno Generator zamków w edytorze}{fig:okno-generatora}{Zrzut okna z \kol{GeneratorWindow}.}

\MiejsceNaRysunek{prefab}{Mapa zapisana jako prefab}{fig:prefab}{Zrzut prefabu z \kol{MapBaker}.}

\MiejsceNaRysunek{aplikacja}{Aplikacja podczas badania}{fig:aplikacja}{Zrzut okna z panelem badań i zamkiem.}

\MiejsceNaRysunek[0.5\linewidth]{atlas}{Atlas materiałów}{fig:atlas}{Tekstura z \kol{MaterialAtlas}.}

\MiejsceNaRysunek{profiler}{Znaczniki w oknie Unity Profiler}{fig:profiler}{Zrzut ze znacznikami \kol{WFC.*}, \cref{tab:znaczniki}.}

\MiejsceNaRysunek{e1-czas}{E1: czas generacji według rozwiązywacza}{fig:e1-czas}{Wykres pudełkowy, kolumna \kol{czas_ms}: pudełko od Q1 do Q3, wąsy do 1,5 rozstępu międzykwartylowego, kropki to pojedyncze mapy.}

\MiejsceNaRysunek[\linewidth]{e2-czas-udane}{E2: czas generacji i odsetek udanych prób według heurystyki}{fig:e2-czas-udane}{Panel a: wykres pudełkowy kolumny \kol{czas_ms}; panel b: \kol{udane_proby_proc} liczone jak w \kol{*_warianty.csv}.}

\begin{figure}
  \centering
  \begin{subfigure}[t]{0.32\linewidth}
    \centering
    \MiejsceNaObraz[\linewidth]{e2-levelfirst}{Zrzut zamku}
    \caption{LevelFirst}
  \end{subfigure}\hfill
  \begin{subfigure}[t]{0.32\linewidth}
    \centering
    \MiejsceNaObraz[\linewidth]{e2-scanline}{Zrzut zamku}
    \caption{Scanline}
  \end{subfigure}\hfill
  \begin{subfigure}[t]{0.32\linewidth}
    \centering
    \MiejsceNaObraz[\linewidth]{e2-minremaining}{Zrzut zamku}
    \caption{MinRemaining}
  \end{subfigure}

  \medskip
  \begin{subfigure}[t]{0.32\linewidth}
    \centering
    \MiejsceNaObraz[\linewidth]{e2-entropy}{Zrzut zamku}
    \caption{Entropy}
  \end{subfigure}\hspace{0.02\linewidth}
  \begin{subfigure}[t]{0.32\linewidth}
    \centering
    \MiejsceNaObraz[\linewidth]{e2-random}{Zrzut zamku}
    \caption{Random}
  \end{subfigure}
  \caption{E2: zamki z tego samego ziarna dla pięciu heurystyk}
  \label{fig:e2-zamki}
\end{figure}

\MiejsceNaRysunek{e3-czas-rozmiar}{E3: czas generacji wobec rozmiaru siatki}{fig:e3-czas-rozmiar}{Mediana i rozstęp międzykwartylowy kolumny \kol{czas_ms} według \kol{sprzecznosc}; obie osie w skali logarytmicznej. Linia limitu 600 s pojawia się, gdy wyniki do niego dochodzą.}

\MiejsceNaRysunek{e4-czas-log}{E4: czas generacji wobec liczby komórek}{fig:e4-czas-log}{Mediana i rozstęp międzykwartylowy kolumny \kol{czas_ms} według \kol{rozwiazywacz}; obie osie w skali logarytmicznej. Nachylenie $a$ (\cref{eq:nachylenie}) skrypt wypisuje w konsoli.}

\MiejsceNaRysunek{e4-pamiec-log}{E4: pamięć wobec liczby komórek}{fig:e4-pamiec-log}{Jak wyżej dla \kol{pamiec_rzeczywista_mb}. Mediana równa 0 nie mieści się na skali logarytmicznej; skrypt wypisuje wtedy, czego brakuje.}

\MiejsceNaRysunek{e5-czas-stany}{E5: czas generacji wobec liczby stanów}{fig:e5-czas-stany}{Mediana i rozstęp międzykwartylowy kolumny \kol{czas_ms} według \kol{zestaw}; stany z \cref{tab:zestawy}.}

\MiejsceNaRysunek[\linewidth]{e6-udane-zabudowa}{E6: odsetek udanych prób i zabudowa według wariantu reguł}{fig:e6-udane-zabudowa}{Kolumny \kol{udane_proby_proc} i \kol{zabudowa_proc} liczone jak w \kol{*_warianty.csv}.}

\MiejsceNaRysunek[\linewidth]{e7-wywolania}{E7: wywołania rysowania wobec rozmiaru siatki}{fig:e7-wywolania}{Mediana i rozstęp międzykwartylowy kolumny \kol{wywolania_rysowania}; panele według \kol{laczenie_rysowania}, linie według \kol{rysowanie}.}

\MiejsceNaRysunek[\linewidth]{e7-klatka}{E7: czas klatki wobec rozmiaru siatki}{fig:e7-klatka}{Jak wyżej dla \kol{czas_klatki_ms}.}

\MiejsceNaRysunek[\linewidth]{e8-restrykcyjnosc}{E8: restrykcyjność reguł a odsetek udanych prób}{fig:e8-restrykcyjnosc}{Wykres punktowy: \kol{gestosc_regul_proc}, \kol{gestosc_wazona_proc} i \kol{kappa} wobec \kol{udane_proby_proc}; pusty znacznik to mapy zdegenerowane.}

\MiejsceNaRysunek{e9-zaostrzanie}{E9: odsetek udanych prób wobec zaostrzenia reguł}{fig:e9-zaostrzanie}{Kolumny \kol{udane_proby_proc} i \kol{zabudowa_proc} według \kol{zaostrzenie_proc}; kropka to jedno losowanie reguł, linia to mediana losowań.}

\MiejsceNaRysunek{propagacja-udzial}{Udział propagacji w usuniętych możliwościach}{fig:propagacja-udzial}{Mapy rozwiązywacza \kol{Bitset} z E4; kolumny \kol{usuniete_propagacja}, \kol{usuniete_decyzje}, \cref{eq:propagacja}.}
